\documentclass[a4paper,11pt,oneside]{article}

% =============================================================================
% DATI DEL FRONTESPIZIO -- da compilare
% =============================================================================
\newcommand{\studente}{Nome Cognome}
\newcommand{\matricola}{0001234567}
\newcommand{\corsodistudi}{Laurea Magistrale in Ingegneria e Scienze Informatiche}
\newcommand{\titoloattivita}{Titolo dell'attività svolta}
\newcommand{\tpv}{TPV} % es. nome/codice dell'attività di tirocinio
\newcommand{\cfu}{12}
\newcommand{\azienda}{Nome Azienda S.p.A.}
\newcommand{\tutoraziendale}{Nome Cognome}
\newcommand{\tutoraccademico}{Prof.\ Nome Cognome}
\newcommand{\periodo}{dal 01/03/2026 al 31/05/2026}
\newcommand{\annoaccademico}{2025/2026}

% Bozza: mostra le indicazioni di compilazione e le note (todonotes).
% Per la versione da consegnare: sostituire con \bozzafalse
\newif\ifbozza
\bozzatrue

% =============================================================================
% PACCHETTI
% =============================================================================
\usepackage{iftex}
\ifPDFTeX
  \usepackage[T1]{fontenc}
\fi
% provide=* carica l'italiano dai file .ini di babel: non serve babel-italian
\usepackage[italian,provide=*]{babel}
\usepackage{lmodern}
\usepackage[varqu,varl,scaled=0.95]{zi4} % Inconsolata come font monospaziato
\usepackage{microtype}
\usepackage[margin=2.5cm,marginparwidth=2cm]{geometry}
\usepackage{graphicx}
\usepackage{booktabs}
\usepackage{tabularx}
\usepackage{xcolor}
\usepackage{listings}
\usepackage[font=small,labelfont=bf]{caption}
\usepackage{url}
\ifbozza
  \usepackage[textsize=small]{todonotes}
\else
  \usepackage[disable]{todonotes}
\fi
\usepackage[hidelinks]{hyperref}
\usepackage[italian,capitalise,noabbrev]{cleveref}
\crefname{listing}{Listato}{Listati}

\hypersetup{
  pdftitle={\titoloattivita},
  pdfauthor={\studente},
}

% Stile del codice (colori ispirati al tema chiaro degli IDE JetBrains)
\definecolor{codebg}{HTML}{F6F8FA}
\definecolor{coderule}{HTML}{4A6FA5}
\definecolor{codekeyword}{HTML}{0033B3}
\definecolor{codestring}{HTML}{067D17}
\definecolor{codecomment}{HTML}{8C8C8C}
\definecolor{codenumber}{HTML}{A0A7B4}
\lstset{
  basicstyle=\ttfamily\small,
  keywordstyle=\color{codekeyword}\bfseries,
  stringstyle=\color{codestring},
  commentstyle=\color{codecomment},
  backgroundcolor=\color{codebg},
  frame=l,
  framerule=2pt,
  rulecolor=\color{coderule},
  framesep=0pt,
  xleftmargin=2.6em,
  framexleftmargin=2.2em,
  framextopmargin=6pt,
  framexbottommargin=6pt,
  numbers=left,
  numbersep=1em,
  numberstyle=\scriptsize\ttfamily\color{codenumber},
  columns=fullflexible,
  keepspaces=true,
  showstringspaces=false,
  upquote=true,
  tabsize=4,
  breaklines=true,
  postbreak=\mbox{\textcolor{codecomment}{$\hookrightarrow$}\space},
  captionpos=b,
  aboveskip=1em,
  belowskip=0.5em,
  % Lettere accentate nel codice (necessario con pdflatex)
  literate=
    {à}{{\`a}}1 {è}{{\`e}}1 {é}{{\'e}}1 {ì}{{\`i}}1 {ò}{{\`o}}1 {ù}{{\`u}}1
    {À}{{\`A}}1 {È}{{\`E}}1 {É}{{\'E}}1 {Ì}{{\`I}}1 {Ò}{{\`O}}1 {Ù}{{\`U}}1,
}
\renewcommand{\lstlistingname}{Listato}
\renewcommand{\lstlistlistingname}{Elenco dei listati}

% Indicazioni per la compilazione: visibili solo in bozza
\newcommand{\guida}[1]{\todo[inline,color=blue!8,bordercolor=blue!40,linecolor=blue!40]{\textbf{Indicazioni:} #1}}
% Note dei tutor allo studente durante la revisione
\newcommand{\notatutor}[1]{\todo[inline,color=orange!20,bordercolor=orange,linecolor=orange]{\textbf{Tutor:} #1}}

% =============================================================================
% DOCUMENTO
% =============================================================================
\begin{document}

\begin{titlepage}
  \centering
  {\large\scshape Alma Mater Studiorum -- Università di Bologna\par}
  \vspace{0.5em}
  {\corsodistudi\par}
  \vspace{0.3em}
  {\small Anno Accademico \annoaccademico\par}

  \vfill

  {\large Relazione sull'attività di \tpv{} (\cfu~CFU)\par}
  \vspace{1.5em}
  {\LARGE\bfseries \titoloattivita\par}
  \vspace{1.5em}
  {\large presso \azienda\par}

  \vfill

  \begin{minipage}[t]{0.48\textwidth}
    \raggedright
    \textbf{Tutor accademico}\\
    \tutoraccademico\\[1em]
    \textbf{Tutor aziendale}\\
    \tutoraziendale
  \end{minipage}%
  \hfill
  \begin{minipage}[t]{0.48\textwidth}
    \raggedleft
    \textbf{Studente}\\
    \studente\\[1em]
    \textbf{Matricola}\\
    \matricola
  \end{minipage}

  \vfill

  {Periodo di svolgimento: \periodo\par}
\end{titlepage}

% =============================================================================
% SOMMARIO
% =============================================================================
\section*{Sommario}
\addcontentsline{toc}{section}{Sommario}

\guida{Massimo una pagina. Riassumere in modo autocontenuto il contesto
(azienda, team, dominio), gli obiettivi dell'attività, il lavoro svolto
personalmente e i principali risultati raggiunti.
Chi legge solo questa pagina deve capire cosa è stato fatto e con quale esito.}

Testo del sommario.

\clearpage

\tableofcontents
\clearpage

% =============================================================================
% INTRODUZIONE
% =============================================================================
\section{Introduzione}
\label{sec:introduzione}

\subsection{Contesto}
\label{sec:contesto}

\guida{Descrivere l'azienda/ente ospitante, il suo dominio di attività
e il progetto o prodotto all'interno del quale si colloca l'attività svolta.}

Testo.

\subsection{Obiettivi}
\label{sec:obiettivi}

\guida{Elencare gli obiettivi dell'attività. Se possibile, renderli verificabili:
nelle conclusioni (\cref{sec:risultati}) si dovrà indicare per ciascuno se e in che misura è stato raggiunto.}

\begin{enumerate}
  \item Primo obiettivo.
  \item Secondo obiettivo.
\end{enumerate}

\subsection{Organizzazione del lavoro}
\label{sec:organizzazione}

\guida{Descrivere il team in cui si è stati inseriti (dimensione, ruoli),
gli stakeholder con cui si è interagito (es.\ colleghi, responsabili, clienti, utenti),
il processo di sviluppo adottato (es.\ Scrum, kanban, ...),
come sono stati fissati gli obiettivi e come è stato monitorato l'avanzamento
(riunioni periodiche, issue tracker, review, ...).}

Testo.

% =============================================================================
% ATTIVITÀ SVOLTE
% =============================================================================
\section{Attività svolte}
\label{sec:attivita}

\guida{Descrivere nel dettaglio ciò che è stato svolto \emph{personalmente},
distinguendolo chiaramente dal lavoro del resto del team.
Scendere nel tecnico: tecnologie, architettura, scelte progettuali e loro motivazioni,
problemi incontrati e soluzioni adottate.}

\subsection{Tecnologie adottate}
\label{sec:tecnologie}

\guida{Elencare le tecnologie (linguaggi, framework, strumenti, servizi)
utilizzate, specificando per ciascuna a cosa è servita.
Le fonti (documentazione, articoli) vanno citate.}

Testo, con eventuali citazioni~\cite{esempio}.

\begin{table}[ht]
  \centering
  \begin{tabularx}{\textwidth}{@{}lX@{}}
    \toprule
    \textbf{Tecnologia} & \textbf{Utilizzo} \\
    \midrule
    Tecnologia A & Descrizione dell'utilizzo. \\
    Tecnologia B & Descrizione dell'utilizzo. \\
    \bottomrule
  \end{tabularx}
  \caption{Tecnologie adottate durante l'attività.}
  \label{tab:tecnologie}
\end{table}

\subsection{Task svolti}
\label{sec:task}

\guida{Descrivere i task svolti, idealmente in ordine cronologico o per area.
Per ciascuno: obiettivo, cosa è stato fatto, eventuali difficoltà e come sono state affrontate.
Usare figure (es.\ diagrammi architetturali, screenshot) e frammenti di codice dove utili.}

\subsubsection{Primo task}

Testo. Si veda ad esempio \cref{fig:architettura} e \cref{lst:esempio}.

\begin{figure}[ht]
  \centering
  \includegraphics[width=0.6\textwidth]{example-image}
  \caption{Architettura del sistema.}
  \label{fig:architettura}
\end{figure}

\begin{lstlisting}[language=Python,caption={Frammento di codice di esempio.},label={lst:esempio}]
# Calcola la media dei valori: è solo un esempio
def media(valori: list[float]) -> float:
    if not valori:
        raise ValueError("lista vuota")
    return sum(valori) / len(valori)
\end{lstlisting}

\subsubsection{Secondo task}

Testo.

% =============================================================================
% CONCLUSIONI
% =============================================================================
\section{Conclusioni e risultati raggiunti}
\label{sec:conclusioni}

\subsection{Risultati raggiunti}
\label{sec:risultati}

\guida{Per ciascuno degli obiettivi fissati in \cref{sec:obiettivi},
indicare se e in che misura è stato raggiunto, con evidenze concrete
(funzionalità rilasciate, metriche, feedback ricevuti, ...).}

Testo.

\subsection{Analisi retrospettiva}
\label{sec:retrospettiva}

\guida{Riflettere criticamente sull'esperienza: cosa ha funzionato e cosa no,
cosa si farebbe diversamente, quali competenze (tecniche e trasversali)
sono state acquisite, quanto il percorso di studi ha preparato all'attività.}

Testo.

\subsection{Sviluppi futuri}
\label{sec:sviluppi}

\guida{Facoltativo: possibili estensioni o prosecuzioni del lavoro svolto.}

Testo.

% Se non si cita alcuna fonte, commentare le due righe seguenti
% (altrimenti compare una bibliografia vuota).
\bibliographystyle{plain}
\bibliography{bibliografia}

\end{document}
